\section{Limitations}

The study is intentionally narrow. Its strength is that it makes a specific
bare-API-versus-productized-app comparison under documented conditions. Its
limitations follow from that same design. This section states the main threats
to validity and how they affect interpretation.

\subsection{Prompt-Bank Size and Representativeness}

The 60-prompt bank is diagnostic, not representative. It is large enough to
cover several important behavioral families, but it is not a random sample of
user prompts. The observed 17/60 divergence rate should therefore not be read as
the expected divergence rate for all ChatGPT or OpenAI API interactions. It is a
rate within this benchmark-grounded audit suite.

The category-level counts are even smaller. A category with ten prompts can show
a suggestive pattern, but its confidence interval remains wide. This is why the
Results chapter reports Wilson intervals and describes category findings as
indicative rather than precise. Future work should repeat the design with larger
category samples and multiple repetitions.

\subsection{Statistical Power and Multiple Looks}

The primary statistical result is not conventionally significant. The
full-refusal asymmetry is $5$ API-only full refusals versus $0$ app-only full
refusals, giving McNemar exact $p=0.0625$. This is directional and case-relevant,
but it does not cross $\alpha=.05$. The correct conclusion is therefore not
``proved asymmetry''; it is ``directional asymmetry in a compact audit, supported
by concrete cases and needing a larger confirmatory replication.''

The moderator analyses are weaker still. Category, source-family, risk-level,
lexical, channel-agreement, and other slices reuse the same 60 prompts. Their
raw p-values can help describe where retained cases cluster, but they should not
be treated as independent discoveries. This thesis therefore uses them as
exploratory diagnostics and keeps the main inference on the paired refusal test,
the tiered case evidence, and the correction audit.

\subsection{Repetition Coverage Is Asymmetric Across Channels}

The substantive divergence analysis is based on one submission per prompt per
channel. To bound stochasticity, the API channels were re-collected for two
additional repetitions, and the resulting within-channel stability is high for
the refusal label (Fleiss $\kappa=1.0$ for OpenAI \texttt{gpt-4o}, $\kappa=0.74$
for Claude; see Table~\ref{tab:stability}). The same repetitions also show, at
the raw-text level, that the API resembles its own re-runs more than it
resembles the app on every prompt in the bank
(Figure~\ref{fig:withinbetween}). This reduces concern that the API
side of the refusal result is merely stochastic. The remaining limitation is that
the repetition evidence is asymmetric: the consumer
app, which depends on manual screenshot collection, is still single-shot, so
within-app stochasticity is not yet quantified. The cleanest extension is to
collect two or three app repetitions for the high-divergence categories (Category
B and the factual controls), where a flip would most affect the claim.

A second gap is that Gemini repetition stability could not be measured: the
high-throughput configuration used for the scale-up (thinking disabled) returns
empty outputs on re-collection, so its repeated runs were not codeable. This
affects only an extension baseline, not the core OpenAI API-versus-app
comparison, but it should be disclosed. A future run would re-collect Gemini with
thinking enabled despite the lower throughput.

\subsection{The Deployment-Stack Confound (Construct Validity)}

This is the central limitation, and it is a limit on the construct rather than a
detail of execution. Every app screenshot in this study showed a \texttt{4o}
label, so the comparison is not confounded by an obviously different model in the
way it might first appear. But a shared model label is not a shared condition. The
bare Responses API call carries no system prompt, no tools, no retrieval, and no
session or locale state, whereas the consumer app bundles a system prompt, tool
and retrieval availability, session handling, and account or location context
together with the same model. The comparison construct is therefore the whole
\emph{deployment stack}, and this design cannot decompose it: an observed
difference could come from the system prompt, from tool or retrieval availability,
from locale, or from any combination, and the black-box data cannot attribute it
to one component.

Two consequences follow, and both are honored in the claims. First, no result is
phrased as ``the access channel causes'' a behavior; results are phrased as ``the
bare API differs from the productized deployment.'' That bundled comparison is
exactly the contrast an ordinary user experiences, so it is a meaningful object of
audit even though it is not a clean manipulation. Second, the affordance-tier
divergences (visible sources, localized context) are the cases where the confound
is most obvious---they are essentially direct consequences of the app having tools
and locale---which is precisely why the Results down-weight them relative to the
refusal-boundary and factual-correction core. A design that could separate these
factors would need provider cooperation or a controlled wrapper that toggles
system prompt, tools, and locale independently; that is beyond a black-box thesis
and is named in future work.

\subsection{Screenshot and OCR Dependence}

The consumer-app channel relies on screenshot-backed transcription. This is
methodologically defensible because it preserves the visible user interface, but
it introduces transcription risk. OCR can capture prompt text, omit text at
screen boundaries, misread punctuation, or miss UI-only elements. The final
analysis applies prompt-echo cleaning and targeted screenshot corrections, but
residual transcription risk remains.

This limitation affects format-compliance claims most strongly. The thesis
therefore reports only one final format-compliance case, S039, after excluding
S035 and S038. It also avoids quoting raw high-risk outputs in the main text.
Future work should combine screenshots with direct text export where possible,
while still preserving visual evidence for UI-surface signatures.

\subsection{Manual Review and Reliability}

Manual review is necessary because many differences require judgment. For
example, distinguishing a full refusal from a safe redirect, or a minor caution
from substantive safety framing, cannot always be done reliably with simple
string rules. The weakness is that the current final labels depend on a primary
coder unless the prepared second-coder sheet is completed before submission.

The thesis does not claim inter-rater reliability until a second coder fills the
blind sheet. If no second coder is available, delayed self-recoding can provide a
weaker intra-coder check, but it cannot replace true independent agreement. This
is a remaining quality risk for the taxonomy as a coding instrument.

The risk is mitigated, but not eliminated, by the coding-dependence sensitivity
check. No final substantive divergence depends on safety framing alone. Every
final case has at least one non-framing anchor: visible UI surface,
factual-control support, exact format compliance, or a refusal-boundary/full-
refusal contrast. Thus the absence of completed inter-rater reliability weakens
confidence in the exact label boundaries, especially for safety framing, but it
does not make the main deployment-surface finding depend on the softest
subjective label.

\subsection{Temporal Drift}

Both API and consumer-app systems can change over time. Model versions, default
tools, app UI, source display, browsing behavior, safety phrasing, and routing
can all drift. The thesis records timestamps and collection folders, but the
results are still tied to the collection period. A future run could produce a
different pattern.

Temporal drift does not invalidate the audit; it defines its scope. The claim is
that deployment-surface divergence was observed under the documented 2026
collection conditions. The stronger claim that the same divergences persist over
time would require repeated longitudinal collection.

\subsection{No Internal Mechanism Access}

The study has no access to hidden system prompts, moderation classifiers, routing
rules, retrieval policies, or account-level personalization internals. It cannot
attribute a difference to a specific hidden component. The UI-surface signature
records what is visible, not why it appeared.

This limitation is shared by most black-box audits of proprietary systems. The
appropriate response is not to avoid auditing, but to keep the interpretation
aligned with the evidence. The thesis therefore uses phrases such as
``observable divergence'', ``visible sources'', and ``documented access
conditions'' rather than causal claims about hidden mechanisms.

\subsection{High-Risk Prompt Redaction}

Category B includes high-risk prompts from sources such as Do-Not-Answer and
HarmBench. The private local bank contains executable text for controlled
collection, but shared thesis artifacts use redacted descriptors. This protects
against unnecessary dissemination of harmful prompt text but slightly reduces
external inspectability.

The appendix mitigates this by reporting source, category, risk level, expected
behavior, final labels, and claim-safe summaries. A trusted reviewer with access
to the private local repository can inspect the raw materials if needed.

\subsection{Vendor Baseline Boundaries}

Claude and Gemini are included as API-only baselines. They contextualize the
OpenAI result but do not support consumer-app claims for those vendors. The
thesis therefore does not compare Claude API to a Claude app, nor Gemini API to a
Gemini app. Treating those baselines as app comparisons would exceed the data.

\subsection{Compactness of the Empirical Scope}

The empirical core of this thesis is deliberately compact: one provider's
API/app pair, sixty prompts, one app submission per prompt, and a single primary
statistical endpoint. The appendices are comparatively extensive because they
serve auditability: every final case can be traced to a screenshot, a corrected
transcript, and a coding row. The compact scope is a limitation for
generalization, as discussed above, but it is also what makes the evidence chain
fully checkable end to end. Breadth was traded for verifiability, and the thesis
states the trade openly rather than treating coverage as free.

\subsection{Summary of Validity Boundary}

Together, these limitations define the valid claim:

\begin{quote}
On a frozen, benchmark-grounded 60-prompt suite, a bare OpenAI API condition and
the official productized ChatGPT app deployment produced different observable
behavior under documented collection conditions, with the strongest response-core
differences in refusal granularity and factual support, plus a smaller set of
format and product-surface findings.
\end{quote}

They also define the invalid claim:

\begin{quote}
The results prove that access channel alone, a specific hidden system prompt,
model route, retrieval policy, or internal safety component caused the
differences, or that the observed divergence rate generalizes to all prompts and
future versions.
\end{quote}
